% Exercise 15.10
\begin{itemize}
    \item[(a)] In Figure~\ref{fig15_10_i}, we have a plot of position as a function of time
        for the Lennard-Jones force with $r_e = D_e = m = 1$, starting from 
        $r = 1.01r_e\,\text{m}$ and $v = 0\,\text{m/s}$. The oscillatory behaviour is
        clear and we can estimate the period at approximately $0.75\,\text{s}$ from reading
        the graph. If we calculate the period using 
        \begin{equation*}
            \begin{split}
                T &= \frac{2\pi}{\omega} \\
                  &= 2\pi r_e \sqrt{\frac{m}{72D_e}} \\
                  &= 2\pi \sqrt{\frac{1}{72}} \\
                  &\approx 0.7405\,\text{s},
            \end{split}
        \end{equation*}
        we can see that our graph is fairly accurate.
    \item[(b)] In Figure~\ref{fig15_10_ii}, we have the same graph for initial position 
        $r = 5r_e$. The graph here does not oscillate, so we cannot evaluate its period.
\end{itemize}

\begin{figure}[!ht]
    \includegraphics[width=\textwidth, center]{../plots/15_10_i.pdf}
    \caption{Position versus time for the Lennard-Jones force with initial position $r = 1.01r_e$ \label{fig15_10_i}} 
\end{figure}

\begin{figure}[!ht]
    \includegraphics[width=\textwidth, center]{../plots/15_10_ii.pdf}
    \caption{Position versus time for the Lennard-Jones force with initial position $r = 5r_e$ \label{fig15_10_ii}} 
\end{figure}
